\chapter{Introduction}
\label{ch:introduction}

Rigid glass beads in nearly incompressible PDMS produce a nonlinear effective
response. Direct finite-element squared (FE\(^2\)) analysis solves an RVE at
every macroscopic integration point and load step, which is costly for repeated
component simulations. An offline constitutive surrogate can reduce this cost
if it supplies a consistent stress and tangent to the finite-element solver.

The primary programming objective is therefore to use a PyTorch pICNN inside
an NGSolve nonlinear solve. Phase~1 trained the pICNN against an analytical
compressible neo-Hookean reference, computed stress and tangent by automatic
differentiation, transferred them through an integration-point bridge, and
verified the resulting equilibrium solve. Only after this proof-of-concept
gate passed did the project proceed to periodic RVE simulations, mixed
near-incompressible matrix mechanics, exact rigid-particle kinematics,
reaction-based homogenization, and a volume-fraction-conditioned surrogate.
This staged work tests the FE material interface before committing to
extensive RVE data generation.

\section{Objectives and scope}

The registered objective was to implement a thermodynamically consistent
neural hyperelastic surrogate for multiscale homogenization of PDMS--glass
bead composites. The work was organized in two stages:
\begin{enumerate}
  \item establish and verify the PyTorch--NGSolve material interface using an
        analytical neo-Hookean reference problem; and
  \item construct the periodic particle RVE, generate homogenized data, train
        the context-dependent surrogate, and deploy it in a macroscopic
        nonlinear finite-element problem.
\end{enumerate}

The implementation replaced the proposed rigid-stiffness penalty with exact
rigid kinematics and the proposed finite-difference homogenized tangent with
symbolic linearization and Schur condensation. Free particle rotation and the
periodic translation gauge were also verified. First-order macroscopic runs
with the Phase~6/7 checkpoint reached 5\% stretch at three mesh sizes, while
the selected shear-adapted checkpoint completed a matched run at one mesh.
Higher-order stability, quantitative response convergence, and experimental
validation remain open.

\section{Report organization}

\Cref{ch:theory} summarizes the physical and mathematical foundations.
\Cref{ch:provenance,ch:implementation,ch:architecture} document software
provenance, numerical implementation, and program structure.
\Cref{ch:results,ch:verification} present results and verification evidence.
\Cref{ch:performance,ch:manual} cover performance and reproducibility, and
\cref{ch:conclusion} concludes the report.
